\section{Appendix --- Source Code}

The complete source code for Laboratory Work 3.2 is organized below by architectural layer. All source files are available in the project GitHub repository.

\textbf{GitHub Repository:} \url{https://github.com/mcittkmims/es-labs}

% ══════════════════════════════════════════════════════════════════════════
% APPLICATION LAYER: ENTRY POINT
% ══════════════════════════════════════════════════════════════════════════

\subsection{Application Layer: Lab 3.2 Entry Point}

\subsubsection{lab3\_2\_main.h}
\begin{lstlisting}[language=C++, caption=lab3\_2\_main.h --- Lab 3.2 Entry Point Interface, label=lst:appendix_lab3_2_main_h]
% [Full source code - see GitHub repository]
% File: labs/lab/lab3_2/lab3_2_main.h
% Declares setup and loop functions for the dual-sensor temperature
% monitoring system with signal conditioning pipeline.
\end{lstlisting}

\subsubsection{lab3\_2\_main.cpp}
\begin{lstlisting}[language=C++, caption=lab3\_2\_main.cpp --- Lab 3.2 Entry Point Implementation, label=lst:appendix_lab3_2_main_cpp]
% [Full source code - see GitHub repository]
% File: labs/lab/lab3_2/lab3_2_main.cpp
% Initializes STDIO, creates synchronization primitives, spawns
% three FreeRTOS tasks (acquisition, conditioning, display),
% and starts the scheduler.
\end{lstlisting}

% ══════════════════════════════════════════════════════════════════════════
% TASK LAYER
% ══════════════════════════════════════════════════════════════════════════

\subsection{Task Layer: Sensor Acquisition}

\subsubsection{task\_acquisition.h}
\begin{lstlisting}[language=C++, caption=task\_acquisition.h --- Sensor Acquisition Task Interface, label=lst:appendix_task_acquisition_h]
% [Full source code - see GitHub repository]
% File: labs/lab/lab3_2/task_acquisition.h
% Declares the vTaskAcquisition() function for periodic sensor reading.
\end{lstlisting}

\subsubsection{task\_acquisition.cpp}
\begin{lstlisting}[language=C++, caption=task\_acquisition.cpp --- Sensor Acquisition Task Implementation, label=lst:appendix_task_acquisition_cpp]
% [Full source code - see GitHub repository]
% File: labs/lab/lab3_2/task_acquisition.cpp
% Periodically reads NTC thermistor (ADC) and DS18B20 (OneWire) at
% 50 ms intervals using vTaskDelayUntil(). Writes raw readings to
% shared g_sensorData under mutex and signals Task 2 via semaphore.
\end{lstlisting}

\subsection{Task Layer: Signal Conditioning}

\subsubsection{task\_conditioning.h}
\begin{lstlisting}[language=C++, caption=task\_conditioning.h --- Signal Conditioning Task Interface, label=lst:appendix_task_conditioning_h]
% [Full source code - see GitHub repository]
% File: labs/lab/lab3_2/task_conditioning.h
% Declares the vTaskConditioning() function for event-driven
% signal conditioning and threshold alerting.
\end{lstlisting}

\subsubsection{task\_conditioning.cpp}
\begin{lstlisting}[language=C++, caption=task\_conditioning.cpp --- Signal Conditioning Task Implementation, label=lst:appendix_task_conditioning_cpp]
% [Full source code - see GitHub repository]
% File: labs/lab/lab3_2/task_conditioning.cpp
% Applies three-stage conditioning pipeline (saturate -> median -> EWMA)
% to both sensor channels, feeds conditioned values to ThresholdAlert
% FSM, and controls LED indicators based on alert states.
\end{lstlisting}

\subsection{Task Layer: Display and Reporting}

\subsubsection{task\_display.h}
\begin{lstlisting}[language=C++, caption=task\_display.h --- Display and Reporting Task Interface, label=lst:appendix_task_display_h]
% [Full source code - see GitHub repository]
% File: labs/lab/lab3_2/task_display.h
% Declares the vTaskDisplay() function for periodic LCD and
% STDIO report updates.
\end{lstlisting}

\subsubsection{task\_display.cpp}
\begin{lstlisting}[language=C++, caption=task\_display.cpp --- Display and Reporting Task Implementation, label=lst:appendix_task_display_cpp]
% [Full source code - see GitHub repository]
% File: labs/lab/lab3_2/task_display.cpp
% Updates LCD with conditioned EWMA values and alert status (Page 0)
% and conditioning configuration (Page 1). Prints structured STDIO
% reports every 2 seconds showing raw, median, and EWMA values for
% both sensors along with alert states and system statistics.
\end{lstlisting}

% ══════════════════════════════════════════════════════════════════════════
% SERVICE LAYER
% ══════════════════════════════════════════════════════════════════════════

\subsection{Service Layer: Shared Data and Synchronization}

\subsubsection{sensor\_data.h}
\begin{lstlisting}[language=C++, caption=sensor\_data.h --- Shared Data Structures and Pin Mapping, label=lst:appendix_sensor_data_h]
% [Full source code - see GitHub repository]
% File: labs/lab/lab3_2/sensor_data.h
% Declares SensorReadings_t (with raw + conditioned fields),
% AlertStatus_t, FreeRTOS synchronization handles, hardware pin
% mapping, NTC/DS18B20 parameters, signal conditioning parameters
% (median window, EWMA alpha, saturation bounds), threshold alert
% parameters, and FreeRTOS task configuration constants.
\end{lstlisting}

\subsubsection{sensor\_data.cpp}
\begin{lstlisting}[language=C++, caption=sensor\_data.cpp --- Shared State Definitions, label=lst:appendix_sensor_data_cpp]
% [Full source code - see GitHub repository]
% File: labs/lab/lab3_2/sensor_data.cpp
% Defines g_sensorData, g_alertData, xNewReadingSemaphore,
% xSensorMutex, and the sensorDataInit() function that creates
% the binary semaphore and mutex.
\end{lstlisting}

\subsection{Service Layer: SignalConditioner}

\subsubsection{SignalConditioner.h}
\begin{lstlisting}[language=C++, caption=SignalConditioner.h --- Signal Conditioning Pipeline Interface, label=lst:appendix_signal_conditioner_h]
% [Full source code - see GitHub repository]
% File: labs/lib/SignalConditioner/SignalConditioner.h
% Declares the SignalConditioner class implementing a three-stage
% pipeline: saturation (clamping), median filter (sliding window
% with insertion sort), and EWMA (first-order IIR filter).
% All intermediate values accessible via getters for diagnostics.
% Maximum window size: 9 samples (fixed array, no heap allocation).
\end{lstlisting}

\subsubsection{SignalConditioner.cpp}
\begin{lstlisting}[language=C++, caption=SignalConditioner.cpp --- Signal Conditioning Pipeline Implementation, label=lst:appendix_signal_conditioner_cpp]
% [Full source code - see GitHub repository]
% File: labs/lib/SignalConditioner/SignalConditioner.cpp
% Implements process(), reset(), saturate(), computeMedian(),
% and all getter methods. The median computation uses insertion
% sort on a local copy of the circular buffer contents.
% EWMA initializes to the first sample on first call.
\end{lstlisting}

\subsection{Service Layer: ThresholdAlert}

\subsubsection{ThresholdAlert.h}
\begin{lstlisting}[language=C++, caption=ThresholdAlert.h --- Hysteresis-Based Threshold Alert Interface, label=lst:appendix_threshold_alert_h]
% [Full source code - see GitHub repository]
% File: labs/lib/ThresholdAlert/ThresholdAlert.h
% Declares the ThresholdAlert class implementing a 4-state FSM:
% NORMAL -> DEBOUNCE_HIGH -> ALERT_ACTIVE -> DEBOUNCE_LOW -> NORMAL
% with configurable high/low thresholds and debounce count.
\end{lstlisting}

\subsubsection{ThresholdAlert.cpp}
\begin{lstlisting}[language=C++, caption=ThresholdAlert.cpp --- Hysteresis-Based Threshold Alert Implementation, label=lst:appendix_threshold_alert_cpp]
% [Full source code - see GitHub repository]
% File: labs/lib/ThresholdAlert/ThresholdAlert.cpp
% Implements the 4-state FSM for threshold detection with
% hysteresis and debounce filtering. Each transition requires
% N consecutive confirmations before state change.
\end{lstlisting}

% ══════════════════════════════════════════════════════════════════════════
% DRIVER LAYER
% ══════════════════════════════════════════════════════════════════════════

\subsection{Driver Layer: AnalogTempSensor}

\subsubsection{AnalogTempSensor.h}
\begin{lstlisting}[language=C++, caption=AnalogTempSensor.h --- NTC Thermistor Driver Interface, label=lst:appendix_analog_temp_h]
% [Full source code - see GitHub repository]
% File: labs/lib/AnalogTempSensor/AnalogTempSensor.h
% Declares the AnalogTempSensor class for reading temperature from
% an NTC thermistor via ADC with Steinhart-Hart Beta equation.
% Circuit: VCC --- [R_series] --- ADC_PIN --- [NTC] --- GND
\end{lstlisting}

\subsubsection{AnalogTempSensor.cpp}
\begin{lstlisting}[language=C++, caption=AnalogTempSensor.cpp --- NTC Thermistor Driver Implementation, label=lst:appendix_analog_temp_cpp]
% [Full source code - see GitHub repository]
% File: labs/lib/AnalogTempSensor/AnalogTempSensor.cpp
% Implements ADC reading, resistance calculation from voltage
% divider equation, and temperature conversion using the
% Steinhart-Hart Beta parameter equation.
\end{lstlisting}

\subsection{Driver Layer: DigitalTempSensor}

\subsubsection{DigitalTempSensor.h}
\begin{lstlisting}[language=C++, caption=DigitalTempSensor.h --- DS18B20 OneWire Driver Interface, label=lst:appendix_digital_temp_h]
% [Full source code - see GitHub repository]
% File: labs/lib/DigitalTempSensor/DigitalTempSensor.h
% Declares the DigitalTempSensor class for reading temperature
% from a DS18B20 sensor via the OneWire protocol.
\end{lstlisting}

\subsubsection{DigitalTempSensor.cpp}
\begin{lstlisting}[language=C++, caption=DigitalTempSensor.cpp --- DS18B20 OneWire Driver Implementation, label=lst:appendix_digital_temp_cpp]
% [Full source code - see GitHub repository]
% File: labs/lib/DigitalTempSensor/DigitalTempSensor.cpp
% Implements OneWire-based temperature reading with non-blocking
% conversion support, device detection, and error handling for
% disconnected and power-on-reset conditions.
\end{lstlisting}

\subsection{Driver Layer: LcdDisplay}

\subsubsection{LcdDisplay.h}
\begin{lstlisting}[language=C++, caption=LcdDisplay.h --- LCD Display Driver Interface, label=lst:appendix_lcd_display_h]
% [Full source code - see GitHub repository]
% File: labs/lib/LcdDisplay/LcdDisplay.h
% Declares the LcdDisplay class wrapping the LiquidCrystal_I2C
% library with convenience methods for two-line updates.
\end{lstlisting}

\subsubsection{LcdDisplay.cpp}
\begin{lstlisting}[language=C++, caption=LcdDisplay.cpp --- LCD Display Driver Implementation, label=lst:appendix_lcd_display_cpp]
% [Full source code - see GitHub repository]
% File: labs/lib/LcdDisplay/LcdDisplay.cpp
% Implements LCD initialization, backlight control, and
% showTwoLines() method for atomic two-line display updates.
\end{lstlisting}

\subsection{Driver Layer: Led}

\subsubsection{Led.h}
\begin{lstlisting}[language=C++, caption=Led.h --- LED Driver Interface, label=lst:appendix_led_h]
% [Full source code - see GitHub repository]
% File: labs/lib/Led/Led.h
% Declares the Led class for controlling a single LED with
% init(), turnOn(), turnOff(), and toggle() methods.
\end{lstlisting}

\subsubsection{Led.cpp}
\begin{lstlisting}[language=C++, caption=Led.cpp --- LED Driver Implementation, label=lst:appendix_led_cpp]
% [Full source code - see GitHub repository]
% File: labs/lib/Led/Led.cpp
% Implements LED control using digitalWrite() with internal
% state tracking for toggle functionality.
\end{lstlisting}

\subsection{Driver Layer: StdioSerial}

\subsubsection{StdioSerial.h}
\begin{lstlisting}[language=C++, caption=StdioSerial.h --- STDIO Serial Redirect Interface, label=lst:appendix_stdio_serial_h]
% [Full source code - see GitHub repository]
% File: labs/lib/StdioSerial/StdioSerial.h
% Declares stdioSerialInit() for redirecting C stdio (printf)
% to the Arduino Serial interface.
\end{lstlisting}

\subsubsection{StdioSerial.cpp}
\begin{lstlisting}[language=C++, caption=StdioSerial.cpp --- STDIO Serial Redirect Implementation, label=lst:appendix_stdio_serial_cpp]
% [Full source code - see GitHub repository]
% File: labs/lib/StdioSerial/StdioSerial.cpp
% Implements STDIO redirection by assigning a custom putchar
% function to the AVR FILE stream, enabling printf() output
% over the hardware UART.
\end{lstlisting}
